\documentclass[a4paper,12pt]{article}
\usepackage[UTF8]{ctex}
\usepackage[backend=bibtex]{biblatex}
\usepackage{graphicx}
\usepackage{abstract}
\usepackage{geometry}
\usepackage{array}
\usepackage{datetime2}
\usepackage{ulem}  % 提供 \uline
\geometry{left=2.5cm, right=2.5cm, top=3cm, bottom=3cm} % 设定页边距
\ctexset{abstractname = {}} % 取消默认摘要标题
\ctexset{contentsname = {}} % 取消默认目录标题

% !编译配方：xelatex -> bibtex -> xelatex*2
% !请务必选择该配方，否则交叉编译目录和引用文献需要编译两次

% 自定义命令：固定长度下划线并居中显示内容
\newcommand{\fixuline}[2][6cm]{\uline{\makebox[#1][c]{#2}}}

% !声明论文主要信息
\newcommand{\LessonName}{\LaTeX 在论文排版中的应用}
\newcommand{\PaperType}{随便什么课程的论文}
\newcommand{\MainTitle}{这是第一个 \LaTeX 模板}
\newcommand{\SubTitle}{——不知道叫什么副标题}
\newcommand{\CollegeName}{材料科学与工程学院}
\newcommand{\AuthorName}{\LaTeX 领域大神}
\newcommand{\ID}{1145141919810}
\newcommand{\Date}{\today} % 如需固定写死，可替换为yyyy-mm-dd

\begin{document}

% --------封面--------
% 插入 logo，宽度设置为 \textwidth
\thispagestyle{empty}
\begin{center}
	\includegraphics[width=\textwidth]{logo.pdf}
\end{center}

\vspace{1cm}

% 课程名称
\begin{center}
	{\Huge \LessonName}\\
	\vspace{0.5cm}
	{\Huge \PaperType}
\end{center}

\vfill

% 主标题
\begin{center}
	{\huge \textbf{\MainTitle}}
\end{center}

% 副标题(不需要请自行注释掉)
\begin{center}
	{\LARGE \SubTitle}
\end{center}

\vfill

% 信息表格，放入 minipage 控制总宽度，避免偏移
\begin{center}
	\begin{minipage}{10cm}
		\renewcommand{\arraystretch}{1.8}
		\begin{tabular}{@{}>{\raggedleft}p{2.5cm}@{\hspace{-0.2cm}}>{\centering\arraybackslash}p{6.7cm}@{}}
			学院： & \fixuline{\CollegeName} \\
			姓名： & \fixuline{\AuthorName} \\
			学号： & \fixuline{\ID} \\
		\end{tabular}
	\end{minipage}
\end{center}

\vfill

% 日期
\begin{center}
	{\large \Date}
\end{center}

% --------摘要--------
\newpage
\pagenumbering{Roman} % 开始罗马数字页码
\begin{abstract}
    \thispagestyle{plain}
    \begin{center}
        {\bfseries\Large 摘\hspace{2em}要} \\[1em]
    \end{center}

    \begin{description}
        \item[目的] 这里是目的
        \item[方法] 这里是方法
        \item[结果] 这里是结果
        \item[结论] 这里是结论
    \end{description}
    
    \vspace{0.5em}
    {\noindent\bfseries 关键词：} 关键词1；关键词2；关键词3
\end{abstract}

% --------目录--------
\newpage
\thispagestyle{empty}
\begin{center}
	{\bfseries\Large 目\hspace{2em}录}
	\tableofcontents % 生成目录
\end{center} 
\pagenumbering{arabic} % 设置页码格式是阿拉伯数字
\setcounter{page}{0}

% --------引言--------
\newpage
\section*{引言} % 星号表示不编号
\addcontentsline{toc}{section}{引言} % 手动添加目录项
这里是引言

% --------正文--------
\section{概述1}
这里是概述1
\subsection{概述1的概述1}
这里是概述1的概述1
\subsection{概述1的概述2}
这里是概述1的概述2

\section{概述2}
这里是概述2
\subsection{概述2的概述1}
这里是概述2的概述1

% --------结语--------
\section*{结语} % 带空格的标题
\addcontentsline{toc}{section}{结语} % 目录项

% --------参考文献--------
\section*{参考文献}
\addcontentsline{toc}{section}{参考文献}
\begin{enumerate}
	\item 参考文献1
	\item 参考文献2
	\item 参考文献3
\end{enumerate}


\end{document}